\chapter{UML Design}

%% ─────────────────────────────────────────────
%% NOTE: This chapter was reconstructed from the thesis PDF. The six UML
%% diagrams (Figures 5.1–5.6) are existing image files — update each
%% \includegraphics path below to point at your actual figure files. The
%% diagrams themselves still depict the old "VAE / Anomaly Detection" stage;
%% they should be regenerated to show the supervised "Normal-vs-Disease gate"
%% so they match the prose and the Implementation/Results chapters.
%% ─────────────────────────────────────────────

\section{Introduction}

This chapter presents the Unified Modeling Language (UML) design adopted for the
proposed mobile application for skin disease detection. UML is used as a visual
modelling language to describe the system and the interactions between its
components in a well-structured form. By using UML diagrams, the system design is
expressed in a way that improves understanding, communication, and
maintainability throughout the software development lifecycle. The design phase
models the essential components of the system, including user interactions,
application workflows, data flow, and the integration of the artificial
intelligence modules. These diagrams confirm that the functional and
non-functional requirements are realised within a robust and scalable software
architecture. The verification and validation of the system---its functional
correctness, performance, and the reliability of its AI predictions and
explanations---are reported separately in the Testing and Evaluation chapter.

\section{Detailed Use Case Diagrams}

\subsection{Skin Analysis}

\begin{figure}[H]
    \centering
    \includegraphics[width=0.85\linewidth]{Figures/usecase_skin_analysis.png}
    \caption{Skin analysis use case diagram.}
    \label{fig:usecase_skin}
\end{figure}

\begin{itemize}
    \item \textbf{Actors:} \emph{General User} — a patient or member of the public; \emph{General Practitioner} — a doctor or frontline clinician.
    \item \textbf{(Primary) Perform Skin Scan:} the main function that both actors initiate.
    \item \textbf{(Extend) Capture New Image:} an optional way to provide an image.
    \item \textbf{(Extend) Upload From Gallery:} a second optional way to provide an image.
    \item \textbf{(Include) Run Normal-vs-Disease Gate:} the system runs the supervised gate to decide whether the image contains a skin condition before any classification.
    \item \textbf{(Include) Run Disease Classification (CNN):} the system runs the B2+B3 CNN ensemble to produce a specific diagnosis (Tinea, Eczema, or Acne).
    \item \textbf{(Include) Generate XAI Heatmap (Score-CAM):} a mandatory final step when a disease is detected — the system provides a visual explanation alongside the diagnosis.
\end{itemize}

\subsection{Administration}

\begin{figure}[H]
    \centering
    \includegraphics[width=0.85\linewidth]{Figures/usecase_admin.png}
    \caption{Administrator use case diagram.}
    \label{fig:usecase_admin}
\end{figure}

\begin{itemize}
    \item \textbf{Actor:} Administrator.
    \item \textbf{Manage User Accounts:} create, delete, or modify user profiles.
    \item \textbf{View System Feedback:} review feedback submitted by general users and clinicians.
    \item \textbf{(Extend) Verify Clinician Credentials:} verify clinician login details.
    \item \textbf{View Performance Dashboard:} access system-wide analytics such as model usage and performance statistics.
\end{itemize}

\subsection{User Account Management}

\begin{figure}[H]
    \centering
    \includegraphics[width=0.85\linewidth]{Figures/usecase_account.png}
    \caption{User account management use case diagram.}
    \label{fig:usecase_account}
\end{figure}

\begin{itemize}
    \item \textbf{Actors:} General User, General Practitioner, and Administrator.
    \item \textbf{Common functions:} both the General User and the General Practitioner can log in / register, manage their own profile, and submit feedback.
    \item \textbf{Administrator functions:} the Administrator additionally manages all user accounts and views system feedback.
\end{itemize}

\subsection{Reports}

\begin{figure}[H]
    \centering
    \includegraphics[width=0.85\linewidth]{Figures/usecase_reports.png}
    \caption{Reports use case diagram.}
    \label{fig:usecase_reports}
\end{figure}

\begin{itemize}
    \item \textbf{Actors:} General User and General Practitioner.
    \item \textbf{View Scan History List:} the base function available to both actors.
    \item \textbf{Select Specific Scan:} optionally select a scan from the history.
    \item \textbf{View Scan Details and Diagnosis:} the system shows the stored diagnosis for the selected scan.
    \item \textbf{Review XAI Heatmap:} the visual explanation is shown together with the diagnosis details.
    \item \textbf{Delete Scan / Export Scan as PDF Report:} two optional actions on a selected scan.
    \item \textbf{Manage Patient Notes:} a function available only to the General Practitioner, allowing professional notes to be attached to a scan.
\end{itemize}

\section{Detailed Sequence Diagram}

Figure~\ref{fig:sequence} shows the detailed flow for analysing a submitted
image. The user provides an image to the \textit{AppUI}, which forwards it to the
\textit{service} layer; the service passes the image to the Normal-vs-Disease
gate.

\begin{itemize}
    \item \textbf{Normal path:} if the gate returns a low $P(\text{disease})$, the service immediately returns ``No Disease Detected'' to the AppUI for display, and the pipeline exits.
    \item \textbf{Disease path:} if the gate returns a high $P(\text{disease})$, the service proceeds to stage two, calling the CNN ensemble to classify the condition (e.g.\ ``Tinea''), then calling the Score-CAM module to generate a matching heatmap.
\end{itemize}

\noindent Finally, the service returns the result (diagnosis and heatmap) to the
AppUI to be displayed to the user.

\begin{figure}[H]
    \centering
    \includegraphics[width=0.9\linewidth]{Figures/sequence_diagram.png}
    \caption{Sequence diagram for the image-analysis workflow.}
    \label{fig:sequence}
\end{figure}

\section{Detailed Class Diagram}

The \texttt{User} class is the parent class for all actors. It holds common
attributes such as \texttt{UserID}, \texttt{Email}, \texttt{Password}, and
\texttt{Role}, and basic operations \texttt{login()}, \texttt{logout()}, and
\texttt{register()}.

\begin{figure}[H]
    \centering
    \includegraphics[width=0.9\linewidth]{Figures/class_diagram.png}
    \caption{Class diagram of the application.}
    \label{fig:class}
\end{figure}

\begin{itemize}
    \item \textbf{Classes inherited from \texttt{User}:}
    \begin{itemize}
        \item \textbf{General User} — can \texttt{performScan()} and \texttt{viewScanHistory()}.
        \item \textbf{Clinician} — holds attributes such as \texttt{MedicalLicenseID} and adds \texttt{viewAllPatientHistory()} and \texttt{managePatientNotes()}.
        \item \textbf{Administrator} — holds privileged operations \texttt{manageUserAccount()}, \texttt{verifyClinician()}, and \texttt{viewSystemFeedback()}.
    \end{itemize}
    \item \textbf{ScanRecord (data class):} holds the key data points — \texttt{UserID}, input image path, diagnosis, confidence score, heatmap image path, and timestamp.
    \item \textbf{HistoryDatabase (service class):} persists data via an SQLite connection and exposes \texttt{saveScan()}, \texttt{getScanHistory()}, and \texttt{deleteScan()}.
    \item \textbf{Relationships:} the General User is associated with the HistoryDatabase, which manages the user's scan history; the HistoryDatabase has a one-to-many relationship with ScanRecord (one database manages many scans).
\end{itemize}

\section{Summary}

This chapter presented the UML design of the proposed Explainable AI powered
mobile application for skin scanning. The use case diagrams captured the
interactions of the general user, clinician, and administrator with the system's
core functions---scanning, classification, explanation, account management, and
reporting. The sequence diagram traced the end-to-end image-analysis workflow
through the Normal-vs-Disease gate, the CNN ensemble, and the Score-CAM
explainability module, and the class diagram defined the data model and its
persistence in the local SQLite store. Together these models provide the
structural and behavioural blueprint that the implementation realises; the
verification of that implementation is reported in the Testing and Evaluation
chapter.
